\documentclass[11pt]{amsart}
\usepackage[localtheorems]{raeez-math-template}

\newtheorem{theorem}{Theorem}[section]
\newtheorem{proposition}[theorem]{Proposition}
\newtheorem{lemma}[theorem]{Lemma}
\newtheorem{definition}[theorem]{Definition}
\newtheorem{remark}[theorem]{Remark}
\newtheorem{example}[theorem]{Example}

\newcommand{\HH}{\mathrm{HH}}
\newcommand{\HC}{\mathrm{HC}}
\newcommand{\ChirHoch}{\mathrm{ChirHoch}}
\newcommand{\Perf}{\mathrm{Perf}}
\newcommand{\Dperf}{D_{\mathrm{perf}}}
\newcommand{\Dfd}{D_{\mathrm{fd}}}
\newcommand{\Db}{D^{b}}
\newcommand{\Coh}{\mathrm{Coh}}
\newcommand{\MF}{\mathrm{MF}}
\newcommand{\Fuk}{\mathrm{Fuk}}
\newcommand{\Rep}{\mathrm{Rep}}
\newcommand{\CoHA}{\mathrm{CoHA}}
\newcommand{\End}{\mathrm{End}}
\newcommand{\Hom}{\mathrm{Hom}}
\newcommand{\Spec}{\mathrm{Spec}}
\newcommand{\Jac}{\mathrm{Jac}}
\newcommand{\NCCR}{\mathrm{NCCR}}
\newcommand{\CY}{\mathrm{CY}}
\newcommand{\Ser}{\mathbf{S}}
\newcommand{\cA}{\mathcal{A}}
\newcommand{\cB}{\mathcal{B}}
\newcommand{\cC}{\mathcal{C}}
\newcommand{\cD}{\mathcal{D}}
\newcommand{\cF}{\mathcal{F}}
\newcommand{\cK}{\mathcal{K}}
\newcommand{\cO}{\mathcal{O}}
\newcommand{\bA}{\mathbb{A}}
\newcommand{\bC}{\mathbb{C}}
\newcommand{\bP}{\mathbb{P}}
\newcommand{\bZ}{\mathbb{Z}}
\newcommand{\fg}{\mathfrak{g}}
\newcommand{\PhiFA}{\Phi^{\mathrm{FA}}}
\newcommand{\SpCh}{\mathrm{SpCh}}
\newcommand{\Zchder}{Z_{\mathrm{ch}}^{\mathrm{der}}}
\newcommand{\op}{\mathrm{op}}

\begin{document}

\title{Oriented Morita Equivalence for Calabi-Yau Threefold Categories}
\author{}
\date{}
\maketitle

\begin{abstract}
The invariant object is an oriented Calabi-Yau dg category
$(\cC,\eta)$, where $\eta$ is a negative cyclic Calabi-Yau class whose
underlying Hochschild class trivialises the inverse dualising bimodule.
Morita equivalence preserves Hochschild and cyclic complexes; it
preserves the oriented datum precisely after the transported orientation
is identified with the chosen one.  Derived equivalence, Fourier-Mukai
equivalence, and CY-linear equivalence are therefore distinct structures.
Geometric categories $\Perf(X)$, Ginzburg dg algebras of quivers with
potential, matrix-factorisation categories, and Fukaya categories supply
different Calabi-Yau sources under different finiteness hypotheses.  In
dimension three the specialised $\Phi_d$ output is $E_1$-chiral; the
$E_2$ structure belongs to the derived chiral centre, Drinfeld centre,
double, or module layer.  The passage from Morita invariance to hCS
quantisation, Hall algebras, Borcherds denominators, and chiral outputs
requires named comparison maps.
\end{abstract}

\tableofcontents

\section{Oriented Calabi-Yau Data}

\begin{definition}[Smooth proper Calabi-Yau dg category]
Let $k$ be a field of characteristic zero.  A small $k$-linear dg
category $\cC$ is smooth if its identity bimodule is perfect as a
$\cC \otimes_k \cC^{\op}$-module, and proper if
$\Hom_\cC(x,y)$ is a perfect $k$-complex for all objects $x,y$.
For smooth proper $\cC$ the inverse dualising bimodule is
\[
  \cC^{!} = \mathrm{RHom}_{\cC^e}(\cC,\cC^e),
  \qquad \cC^e=\cC\otimes_k\cC^{\op}.
\]
A Calabi-Yau structure of dimension $d$ is a negative cyclic class
\[
  \eta\in HC^-_d(\cC)
\]
whose image in $\HH_d(\cC)$ induces an equivalence
$\cC^{!}\simeq \cC[-d]$.  Equivalently, the Serre functor on
$\Dperf(\cC)$ is equipped with a chosen trivialisation
\[
  \Ser_\cC \simeq [d].
\]
The pair $(\cC,\eta)$ is an oriented CY$_d$ dg category.
\end{definition}

The orientation is part of the object.  The underlying dg category
remembers the Serre functor only up to its universal property; the class
$\eta$ chooses the trace and hence the pairing
\[
  \langle-,-\rangle_\eta\colon
  \Hom_\cC(x,y)\otimes\Hom_\cC(y,x)\longrightarrow k[-d].
\]
For $\cC=\Perf(X)$, with $X$ a smooth proper Calabi-Yau $d$-fold, a
holomorphic volume form $\Omega_X$ determines $\eta_X$ through the
Serre residue trace.

\begin{definition}[CY-linear Morita equivalence]
Let $(\cC,\eta)$ and $(\cD,\theta)$ be oriented CY$_d$ dg categories.
A CY-linear Morita equivalence from $(\cC,\eta)$ to $(\cD,\theta)$ is a
Morita equivalence $F\colon\Dperf(\cC)\simeq\Dperf(\cD)$ together with
a specified path
\[
  F_*(\eta)\simeq \theta
  \quad\text{in}\quad HC^-_d(\cD).
\]
\end{definition}

\begin{theorem}[Oriented Morita invariance]
\label{thm:oriented-morita-invariance}
Morita equivalence induces natural isomorphisms
\[
  F_* \colon \HH_\bullet(\cC)\xrightarrow{\sim}\HH_\bullet(\cD),
  \qquad
  F_* \colon HC^-_\bullet(\cC)\xrightarrow{\sim}HC^-_\bullet(\cD).
\]
The transported class $F_*(\eta)$ is a Calabi-Yau structure on $\cD$.
The oriented objects $(\cC,\eta)$ and $(\cD,\theta)$ are equivalent
exactly when $F_*(\eta)$ and $\theta$ are identified in negative cyclic
homology.  Thus the invariant is the oriented Morita class.
\end{theorem}

\begin{proof}
Keller's Morita invariance of Hochschild and cyclic homology gives the
two displayed isomorphisms.  The non-degeneracy condition is expressed
by the bimodule map $\cC^{!}\to\cC[-d]$ induced by the Hochschild image of
$\eta$.  Derived Morita equivalence transports perfect bimodules and
therefore carries this equivalence to
\[
  \cD^{!} \simeq \cD[-d].
\]
This proves that $F_*(\eta)$ is a Calabi-Yau structure on $\cD$.
Equality with the chosen structure $\theta$ is precisely the additional
orientation datum in the definition of CY-linear Morita equivalence.
\end{proof}

\begin{remark}[Serre functor versus orientation]
An exact equivalence between smooth proper categories identifies Serre
functors by uniqueness:
\[
  F\Ser_\cC\simeq\Ser_\cD F.
\]
This canonical compatibility does less than CY-linearity.  CY-linearity
requires compatibility with the chosen trivialisation
$\Ser\simeq[d]$, equivalently with the chosen negative cyclic trace.
For a fixed Calabi-Yau variety $X$, replacing a volume form
$\Omega_X$ by $\lambda\Omega_X$ changes the orientation by
$\lambda\in\bC^\times$.  The identity functor on $\Perf(X)$ is
oriented only after the two scalar choices are identified.
\end{remark}

\begin{proposition}[The Morita layer]
\label{prop:morita-layer}
Let $F$ be a derived Morita equivalence between smooth proper dg
categories $\cC$ and $\cD$, represented by an invertible
$\cC$-$\cD$ bimodule.  Then $F$ gives:
\begin{enumerate}
\item an equivalence of perfect module categories
\[
  \Dperf(\cC)\simeq\Dperf(\cD);
\]
\item a quasi-isomorphism of Hochschild chain complexes and of negative
cyclic complexes;
\item an equivalence of Hochschild cochain $E_2$-algebras
\[
  \mathrm{RHom}_{\cC^e}(\cC,\cC)
  \simeq
  \mathrm{RHom}_{\cD^e}(\cD,\cD),
\]
hence an identification of derived centre data;
\item in a monoidal or $E_1$ presentation, an equivalence of the
corresponding Drinfeld centres.
\end{enumerate}
If $F$ is CY-linear, these identifications preserve the chosen negative
cyclic orientation.  Structures beyond this list require their own
comparison data: a BV kernel for hCS quantisation, a stability and
orientation package for Hall multiplication, an automorphic product for
a Borcherds denominator, and a morphism inside the verified domain of
$\PhiFA_d$ for a chiral output.
\end{proposition}

\begin{proof}
The invertible bimodule identifies the compact module categories and
therefore the categories of perfect bimodules.  Hochschild chains,
cyclic chains, and Hochschild cochains are respectively traces and
endomorphisms of the unit bimodule in this Morita bicategory, so Keller
and To\"en Morita invariance give the displayed identifications.  In a
monoidal presentation, the Drinfeld centre is the endomorphism object of
the identity module category, hence is invariant under equivalence of
the monoidal module category.  CY-linearity adds the path in $HC^-_d$
from the transported orientation to the chosen target orientation,
exactly as in Theorem~\ref{thm:oriented-morita-invariance}.  The
remaining structures use moduli of fields, extension correspondences,
automorphic divisors, or specialisation cycles; those data lie outside
the Morita bicategory until supplied explicitly.
\end{proof}

\section{Fourier-Mukai Equivalence and Flops}

Let $X$ and $Y$ be smooth proper Calabi-Yau $d$-folds over $\bC$, with
volume forms $\Omega_X$ and $\Omega_Y$.  A Fourier-Mukai kernel
$\cK\in\Db\Coh(X\times Y)$ gives an equivalence
\[
  \Phi_\cK\colon\Db\Coh(X)\xrightarrow{\sim}\Db\Coh(Y)
\]
under the usual perfectness and full-faithfulness hypotheses.  The
Hochschild-Kostant-Rosenberg identification and the Mukai vector of
$\cK$ give the induced map on Hochschild homology.

\begin{proposition}[Fourier-Mukai orientation criterion]
\label{prop:fm-orientation-criterion}
The equivalence $\Phi_\cK$ is CY-linear exactly when its Hochschild
action sends the oriented top class of $X$ to the oriented top class of
$Y$:
\[
  \HH_d(\Phi_\cK)([\Omega_X])=[\Omega_Y].
\]
Equivalently, under HKR, the cohomological Fourier-Mukai transform
defined by $v(\cK)=\operatorname{ch}(\cK)\sqrt{\operatorname{td}_{X\times Y}}$
carries the Serre trace associated with $\Omega_X$ to the Serre trace
associated with $\Omega_Y$.
\end{proposition}

\begin{proof}
Caldararu's Mukai-pairing theorem identifies the Hochschild pairing of
$\Perf(X)$ with the Mukai pairing after the Todd correction.  A
Fourier-Mukai equivalence preserves the pairing.  The oriented
Calabi-Yau datum asks for preservation of the chosen generator of the
one-dimensional line $H^0(X,\omega_X)$, transported through the
Hochschild action.  This is exactly the displayed condition.
\end{proof}

\begin{proposition}[Crepant threefold flop]
\label{prop:crepant-flop-oriented}
Let
\[
  f_\pm\colon Y_\pm\longrightarrow X_0
\]
be a projective threefold flop over a Gorenstein base $X_0$ with
crepant resolutions $Y_\pm$.  Fix a reflexive trivialisation of
$\omega_{X_0}$ on the smooth locus, and orient $Y_\pm$ by pulling this
trivialisation back through $f_\pm^!\omega_{X_0}\simeq\omega_{Y_\pm}$.
Then Bridgeland's flop equivalence
\[
  \Db\Coh(Y_+)\simeq\Db\Coh(Y_-)
\]
is CY-linear for these orientations.  The statement concerns the chosen
crepant orientation; changing either volume form by a scalar changes the
oriented equivalence problem.
\end{proposition}

\begin{proof}
Bridgeland constructs the equivalence by a Fourier-Mukai kernel
supported on $Y_+\times_{X_0}Y_-$.  On the common open subset where the
flop is an isomorphism the kernel is the diagonal kernel, and the two
volume forms are the same pullback of the fixed reflexive generator of
$\omega_{X_0}$.  The complement has codimension at least two in the
threefolds, while the canonical sheaves are reflexive.  Hence the
agreement of the canonical trivialisations on the common open subset
extends uniquely.  Proposition~\ref{prop:fm-orientation-criterion}
then gives CY-linearity.
\end{proof}

\begin{example}[The conifold]
Let
\[
  R=\bC[x,y,z,w]/(xy-zw),\qquad X_0=\Spec R,
\]
and let $Y_\pm\to X_0$ be the two small resolutions.  With
$I=(x,z)\subset R$, Van den Bergh's tilting bundles identify both
small resolutions with the same non-commutative crepant resolution
\[
  \Lambda=\End_R(R\oplus I).
\]
The Klebanov-Witten quiver has vertices $0,1$, arrows
$a_1,a_2\colon0\to1$ and $b_1,b_2\colon1\to0$, and potential
\[
  W_{\mathrm{KW}}=a_1b_1a_2b_2-a_1b_2a_2b_1.
\]
Its Jacobi algebra satisfies
\[
  J_{\mathrm{con}}\cong \Lambda,\qquad Z(J_{\mathrm{con}})\cong R.
\]
This is the precise Morita statement in the conifold chart:
\[
  \Db\Coh(Y_\pm)\simeq\Db(\Lambda\text{-}\mathrm{mod})
  \simeq\Db(J_{\mathrm{con}}\text{-}\mathrm{mod}).
\]
The chart is local and non-compact, so its trace is an equivariant or
compact-support trace.  It is a local CY$_3$ model, separate from a
smooth proper global Calabi-Yau datum.
\end{example}

\section{Four Sources of CY3 Categories}

The phrase ``CY3 category'' has several legitimate origins.  The
following list records the source, the orientation, and the finiteness
hypotheses.

\begin{enumerate}
\item \emph{Geometric source.}
If $X$ is a smooth proper Calabi-Yau threefold with a volume form
$\Omega_X$, then $\Perf(X)$ is a smooth proper oriented CY$_3$ dg
category.  The trace is the Serre residue trace associated with
$\Omega_X$.

\item \emph{Ginzburg source.}
For a finite quiver with potential $(Q,W)$, the Ginzburg dg algebra
$\Gamma(Q,W)$ is homologically smooth and bimodule $3$-Calabi-Yau:
\[
  \mathrm{RHom}_{\Gamma^e}(\Gamma,\Gamma^e)\simeq\Gamma[-3].
\]
This is Ginzburg's construction, refined by Keller and Van den Bergh as
a deformed Calabi-Yau completion.  If the Jacobi algebra is
finite-dimensional, $\Dfd(\Gamma(Q,W))$ is a proper CY$_3$ triangulated
category.  For toric local threefolds such as $\bC^3$ and the conifold,
one uses the local or equivariant BPS category with compact-support
pairings.

\item \emph{Matrix-factorisation source.}
Let $S$ be a smooth $n$-dimensional Calabi-Yau variety and
$W\colon S\to\bA^1$ a regular function with proper critical locus.
The dg category $\MF(S,W)$ carries the Kapustin-Li residue trace and,
under the standard isolated-critical or properness hypotheses, is a
proper Calabi-Yau category of dimension $n-2$ in the singularity-category
normalisation.  Thus an affine Landau-Ginzburg potential on $\bA^5$
is the ambient dimension that can supply a CY$_3$ matrix-factorisation
category.  Orlov's LG/CY comparison relates graded matrix
factorisations to $\Db\Coh$ of projective Calabi-Yau hypersurfaces only
under its grading and GIT hypotheses.

\item \emph{Fukaya source.}
For a graded exact or compact symplectic six-manifold with
$c_1=0$, and for a Fukaya category satisfying the usual compactness,
orientation, spin, and transversality hypotheses, the Floer
Poincare-duality trace gives a CY$_3$ $A_\infty$ category.  Wrapped
Fukaya categories require the corresponding locally proper or
compact-object formulation.
\end{enumerate}

These four sources overlap in special examples through homological
mirror symmetry, tilting, Orlov equivalence, or Knorrer periodicity.
They are distinct constructions.  A Ginzburg dg algebra, a
matrix-factorisation category, a Fukaya category, and $\Perf(X)$ carry
different input data, even when a theorem identifies two of them in a
particular example.

\section{Matrix Factorisations and Knorrer Periodicity}

\begin{proposition}[Calabi-Yau dimension of an isolated hypersurface category]
\label{prop:mf-dimension}
Let $W\colon\bA^n\to\bA^1$ have an isolated critical point at the origin.
The matrix-factorisation category $\MF(\bA^n,W)$ has Kapustin-Li
pairing
\[
  \langle f,g\rangle_{\mathrm{KL}}
  =
  \operatorname*{Res}_{0}
  \frac{fg\,dx_1\wedge\cdots\wedge dx_n}
       {\partial_{x_1}W\cdots\partial_{x_n}W}.
\]
In the singularity-category convention its Serre functor is the shift
$[n-2]$.  Hence the Calabi-Yau dimension is $n-2$.
\end{proposition}

\begin{proof}
Dyckerhoff identifies $\MF(\bA^n,W)$ with the dg category generated by
the stabilised residue field and computes its Hochschild chains by the
Jacobian algebra.  Polishchuk-Vaintrob and Shklyarov write the explicit
non-degenerate cyclic cocycle as the Kapustin-Li residue, with higher
corrections in the global case.  For a hypersurface
$R=\bC[x_1,\ldots,x_n]/(W)$ of Krull dimension $n-1$, Buchweitz-Orlov
duality gives Serre shift $[n-2]$ on the stable singularity category,
which is the matrix-factorisation category under Orlov's equivalence.
\end{proof}

\begin{proposition}[Knorrer periodicity]
\label{prop:knorrer}
For $W$ as above,
\[
  \MF(\bA^{n+2},W+uv)\simeq\MF(\bA^n,W).
\]
The equivalence raises the ambient affine dimension by two and shifts
the integral Serre dimension by two.  In the two-periodic
matrix-factorisation category this shift is invisible; in a graded or
oriented comparison it must be recorded.
\end{proposition}

\begin{example}[The potential $xyz$]
For $W=xyz$ on $\bA^3$,
\[
  \partial_xW=yz,\qquad \partial_yW=xz,\qquad \partial_zW=xy,
\]
so
\[
  \operatorname{Crit}(W)=V(xy,xz,yz)
  =
  V(x,y)\cup V(x,z)\cup V(y,z).
\]
The critical locus has dimension one and the Jacobian algebra
$\bC[x,y,z]/(xy,xz,yz)$ is infinite-dimensional.  The properness
hypothesis in Proposition~\ref{prop:mf-dimension} fails.  Even after
replacing $xyz$ by an isolated potential on $\bA^3$, the
matrix-factorisation Calabi-Yau dimension would be $1$, since
$3-2=1$.  Thus $\MF(\bA^3,xyz)$ cannot serve as the CY$_3$ category of
the $\bC^3$ quiver.
\end{example}

\begin{proposition}[The $\bC^3$ quiver datum]
\label{prop:c3-ginzburg-datum}
Let $Q_{\bC^3}$ be the one-vertex quiver with loops $x,y,z$, and let
\[
  W_{\bC^3}=xyz-xzy\in \bC Q_{\bC^3}/[\bC Q_{\bC^3},\bC Q_{\bC^3}].
\]
Then
\[
  \partial_x W_{\bC^3}=yz-zy,\quad
  \partial_y W_{\bC^3}=zx-xz,\quad
  \partial_z W_{\bC^3}=xy-yx,
\]
and
\[
  \Jac(Q_{\bC^3},W_{\bC^3})
  =
  \bC\langle x,y,z\rangle/([x,y],[y,z],[z,x])
  \cong \bC[x,y,z].
\]
The associated Ginzburg dg algebra $\Gamma(Q_{\bC^3},W_{\bC^3})$ is the
3-Calabi-Yau dg algebra attached to the non-commutative $\bC^3$ chart.
The cohomological Hall algebra of this chart is the positive half
\[
  \CoHA(\bC^3)=Y^+(\widehat{\mathfrak{gl}}_1).
\]
The full $\mathcal W_{1+\infty}$ object arises only after the
Drinfeld-double or centre passage.
\end{proposition}

\begin{proof}
The cyclic derivatives of $W_{\bC^3}$ are the three commutators
displayed above, so the Jacobi algebra is the commutative polynomial
algebra in three variables.  Ginzburg's construction adds opposite
arrows and degree-$-2$ loops with differential determined by these
cyclic derivatives, giving a homologically smooth bimodule
$3$-Calabi-Yau dg algebra.  The Schiffmann--Vasserot comparison
identifies the critical cohomological Hall algebra of the
$\bC^3$ chart with the positive half of the affine Yangian of
$\widehat{\mathfrak{gl}}_1$.  The vertex algebra
$\mathcal W_{1+\infty}$ uses the doubled, centred, completed, or
evaluation-realised object.
\end{proof}

\section{Hall, Borcherds, and Chiral Layers}

\begin{definition}[Hall-compatible CY-linear equivalence]
Let $(\cC,\eta)$ and $(\cD,\theta)$ be oriented CY$_3$ dg categories
equipped with stability data, orientation-line square roots, and
vanishing-cycle sheaves on their derived moduli stacks of objects.  A
CY-linear Morita equivalence $F\colon\cC\simeq\cD$ is
Hall-compatible when the induced equivalence of moduli stacks carries
semistable substacks to semistable substacks, identifies the
orientation-line and vanishing-cycle data, and intertwines the
extension correspondence defining Hall multiplication.
\end{definition}

\begin{proposition}[Hall transport]
\label{prop:hall-transport}
A Hall-compatible CY-linear equivalence induces an $E_1$-algebra
quasi-isomorphism
\[
  \CoHA(\cC,\sigma_\cC)
  \simeq
  \CoHA(\cD,\sigma_\cD).
\]
The assertion concerns the Hall positive half attached to the chosen
stability and orientation data.
\end{proposition}

\begin{proof}
The cohomological Hall product is the pull-push operation along the
derived stack of short exact triangles
\[
  \cM_\cC\times\cM_\cC
  \xleftarrow{\;p\;}
  \cM_\cC^{\mathrm{ext}}
  \xrightarrow{\;q\;}
  \cM_\cC,
\]
with coefficients in the chosen vanishing-cycle sheaf and orientation
line.  Hall compatibility identifies this correspondence with the
corresponding one for $\cD$.  Proper base-change and the compatibility
of the coefficient systems transport convolution.  The resulting map is
an equivalence because the moduli-stack and coefficient identifications
come from equivalences.
\end{proof}

\begin{proposition}[Chiral comparison data]
\label{prop:chiral-comparison-data}
Let $F\colon(\cC,\eta)\simeq(\cD,\theta)$ be a CY-linear Morita
equivalence in the verified domain of $\PhiFA_d$.  Suppose
\[
  \PhiFA_d(F)\colon \PhiFA_d(\cC)\longrightarrow\PhiFA_d(\cD)
\]
is constructed and the specialisation data
$(\Sigma_{d-1}^{\cC},C_\cC)$ and $(\Sigma_{d-1}^{\cD},C_\cD)$ are
identified.  Then the specialised outputs are equivalent at their
native operadic level:
\[
  \SpCh_{\Sigma_{d-1}^{\cC},C_\cC}(\PhiFA_d(\cC))
  \simeq
  \SpCh_{\Sigma_{d-1}^{\cD},C_\cD}(\PhiFA_d(\cD)).
\]
For $d\geq3$ this is an $E_1$-chiral equivalence.  A braided $E_2$
comparison is obtained after applying the derived chiral centre,
Drinfeld centre, completed double, or representation-category layer on
the loci where that layer has been constructed.
\end{proposition}

\begin{proof}
The constructed morphism $\PhiFA_d(F)$ lives in the target category of
$E_d$-holomorphic factorisation algebras.  Factorisation homology over
the identified $(d-1)$-cycle and restriction to the identified reference
curve are functorial for such morphisms.  The native level is the
dimension-forced residual level of the two-stage construction:
$E_\infty$ for $d=1$, $E_2$ for $d=2$, and $E_1$ for $d\geq3$.
Centres and doubles are functorial after their own module-category,
Hopf-pairing, or completion data have been supplied.
\end{proof}

\begin{example}[\texorpdfstring{$K3\times E$}{K3 x E} invariants]
Let $X=K3\times E$.  Kunneth gives
\[
  h^{0,\bullet}(X)=(1,1,1,1),
  \qquad
  \kappa_{\mathrm{ch}}(X)=1-1+1-1=0.
\]
The categorical Euler invariant is multiplicative:
\[
  \kappa_{\mathrm{cat}}(X)
  =
  \chi(\cO_{K3})\chi(\cO_E)
  =
  (1+1)(1-1)=0.
\]
The Heisenberg-Mukai specialisation gives
\[
  \kappa_{\mathrm{ch}}^{\mathrm{Heis}}(X)=3.
\]
The Borcherds denominator attached to $\Delta_5$ gives
\[
  \kappa_{\mathrm{BKM}}(\fg_{\Delta_5})=\frac{c_1(0)}{2}
  =\frac{10}{2}=5.
\]
The fibre lattice witness is
\[
  \kappa_{\mathrm{fiber}}=24.
\]
Thus the compact total-space supertraces, the Heisenberg-Mukai
specialisation, the Borcherds weight, and the fibre lattice rank are
four separate constructions.  The Hall-Drinfeld/Borcherds branch and
the K3 self-mirror current branch are separate presentation problems.
\end{example}

\section{The CY-to-Chiral Scope}

The two-stage construction in the main manuscript has the form
\[
  \Phi_d^{(\Sigma_{d-1},C)}
  =
  \SpCh_{\Sigma_{d-1},C}\circ\PhiFA_d,
  \qquad
  \CY_d\text{-}\mathrm{Cat}
  \xrightarrow{\PhiFA_d}
  E_d\text{-}\mathrm{HolFA}(X)
  \xrightarrow{\SpCh_{\Sigma_{d-1},C}}
  E_{n_{\mathrm{native}}(d)}\text{-}\mathrm{HolFA}(C).
\]
Stage $1$ uses the Hochschild cochain complex, the Gerstenhaber bracket
of degree $1-d$, a formality choice, and the Calabi-Yau trace.  At
$d=3$ it is a witnessed construction under the framed Stage-$1$
hypotheses; arbitrary morphism-level functoriality remains a
separate theorem to prove.  Stage $2$ depends on an admissible
specialisation datum $(\Sigma_{d-1},C)$.

\begin{proposition}[Native operadic level]
\label{prop:native-operadic-level-note}
For the specialised chiral algebra
\[
  A_\cC=
  \SpCh_{\Sigma_{d-1},C}\bigl(\PhiFA_d(\cC)\bigr),
\]
the native operadic level is
\[
  n_{\mathrm{native}}(d)=
  \begin{cases}
    \infty, & d=1,\\
    2, & d=2,\\
    1, & d\geq3.
  \end{cases}
\]
In particular, a CY$_3$ input gives an $E_1$-chiral algebra on the
reference curve.  The $E_2$ structure in dimension three is carried by
the centre, double, or module layer, for instance
\[
  \mathcal Z\bigl(\Rep^{E_1}(A_\cC)\bigr)
  \simeq
  \Rep^{E_2}\bigl(\Zchder(A_\cC)\bigr)
\]
on the verified Drinfeld-centre locus.
\end{proposition}

\begin{remark}[Five separated objects]
For an $E_1$-chiral algebra $A$, the objects
\[
  A,\qquad B(A),\qquad A^{i},\qquad A^{!},\qquad \Zchder(A)
\]
have different meanings.  The bar construction is $B(A)$; the identity
$\Omega B(A)\simeq A$ is bar-cobar inversion.  The object $A^{i}$ is the
Koszul dual coalgebra.  The object $A^{!}$ is the Verdier or Koszul dual
algebra, depending on the duality context.  The derived chiral centre
$\Zchder(A)$ is the bulk or centre object carrying the braided layer in
the $d\geq3$ specialisation.
\end{remark}

\section{Morita Scope}

\begin{theorem}[Morita and derived-equivalence scope for CY3 inputs]
\label{thm:morita-derived-scope}
Let $(\cC,\eta)$ be a CY$_3$ input to the verified locus of
$\PhiFA_3$.
\begin{enumerate}
\item The oriented datum is invariant under CY-linear Morita
equivalence, namely a Morita equivalence carrying $\eta$ to the chosen
orientation on the target.
\item A Fourier-Mukai equivalence between smooth proper Calabi-Yau
threefolds identifies the oriented data exactly when its Hochschild
action carries the chosen holomorphic volume form to the chosen
holomorphic volume form.
\item A crepant flop is CY-linear after both sides are oriented by the
same reflexive canonical trivialisation on the common Gorenstein base.
\item Ginzburg dg algebras, matrix-factorisation categories, Fukaya
categories, and categories $\Perf(X)$ are legitimate CY sources under
their own hypotheses.  Equivalences among them are theorems in specific
examples, supplied by their comparison functors.
\item The specialised chiral output of a CY$_3$ source is $E_1$-chiral.
Any $E_2$ enhancement belongs to the derived centre, Drinfeld centre,
double, or representation layer.
\item Morita equivalence preserves perfect module categories,
Hochschild chains, negative cyclic chains, Hochschild cochain
$E_2$-algebras, and centre data under the hypotheses of
Proposition~\ref{prop:morita-layer}.
\item Hall/CoHA comparison requires a Hall-compatible equivalence of
moduli stacks, stability data, orientation lines, vanishing-cycle
sheaves, and extension correspondences.
\item hCS quantisation, Borcherds denominators, and chiral outputs are
identified only after the relevant BV, automorphic, and
two-stage-specialisation comparison maps have been constructed.
\end{enumerate}
\end{theorem}

\begin{proof}
Clause (1) is Theorem~\ref{thm:oriented-morita-invariance}.  Clause (2)
is Proposition~\ref{prop:fm-orientation-criterion}.  Clause (3) is
Proposition~\ref{prop:crepant-flop-oriented}.  Clause (4) follows from
the four-source catalogue above: smooth proper geometry gives Serre
residue traces, Ginzburg dg algebras give bimodule 3-Calabi-Yau
structures and proper CY$_3$ finite-dimensional categories under
Jacobi-finiteness, matrix factorisations give CY dimension $n-2$ under
proper critical-locus hypotheses, and Fukaya categories give cyclic
Floer pairings under the standard grading and compactness hypotheses.
Clause (5) is Proposition~\ref{prop:native-operadic-level-note}, the
dimension-three case of the two-stage CY-to-chiral construction.
Clause (6) is Proposition~\ref{prop:morita-layer}.  Clause (7) is
Proposition~\ref{prop:hall-transport}.  Clause (8) is
Proposition~\ref{prop:chiral-comparison-data} together with the
definition of the Borcherds weight
$\kappa_{\mathrm{BKM}}(\Phi_N)=c_N(0)/2$.
\end{proof}

\section{References}

\begin{itemize}
\item Bondal, A. I., and Kapranov, M. M., \emph{Enhanced triangulated
categories}, Math. USSR Sb. 70 (1991), 93-107.
\item Bridgeland, T., \emph{Flops and derived categories},
Invent. Math. 147 (2002), 613-632.
\item Brav, C., and Dyckerhoff, T., \emph{Relative Calabi-Yau structures},
Compos. Math. 155 (2019), 372-412.
\item Buchweitz, R.-O., \emph{Maximal Cohen-Macaulay modules and Tate
cohomology over Gorenstein rings}, 1986 manuscript.
\item Caldararu, A., \emph{The Mukai pairing, I: the Hochschild pairing},
Adv. Math. 194 (2005), 34-66.
\item Costello, K., \emph{Topological conformal field theories and
Calabi-Yau categories}, Adv. Math. 210 (2007), 165-214.
\item Costello, K., and Gwilliam, O., \emph{Factorization Algebras in
Quantum Field Theory}, Cambridge University Press, 2017 and 2021.
\item Dyckerhoff, T., \emph{Compact generators in categories of matrix
factorizations}, Duke Math. J. 159 (2011), 223-274.
\item Ginzburg, V., \emph{Calabi-Yau algebras}, arXiv:math/0612139, 2006.
\item Keller, B., \emph{On the cyclic homology of exact categories},
J. Pure Appl. Algebra 136 (1999), 1-56.
\item Keller, B., \emph{Deformed Calabi-Yau completions}, with appendix
by M. Van den Bergh, J. Reine Angew. Math. 654 (2011), 125-180.
\item Knorrer, H., \emph{Cohen-Macaulay modules on hypersurface
singularities I}, Invent. Math. 88 (1987), 153-164.
\item Kontsevich, M., and Soibelman, Y., \emph{Cohomological Hall
algebra, exponential Hodge structures and motivic Donaldson-Thomas
invariants}, Commun. Number Theory Phys. 5 (2011), 231-352.
\item Lurie, J., \emph{Higher Algebra}, available from the author's
website.
\item Orlov, D., \emph{Triangulated categories of singularities and
D-branes in Landau-Ginzburg models}, Proc. Steklov Inst. Math. 246
(2004), 227-248.
\item Polishchuk, A., and Vaintrob, A., \emph{Chern characters and
Hirzebruch-Riemann-Roch formula for matrix factorizations},
Duke Math. J. 161 (2012), 1863-1926.
\item Schiffmann, O., and Vasserot, E., \emph{Cherednik algebras,
W-algebras and the equivariant cohomology of the moduli space of
instantons on the plane}, Publ. Math. IHES 118 (2013), 213-342.
\item Seidel, P., \emph{Fukaya Categories and Picard-Lefschetz Theory},
European Mathematical Society, 2008.
\item Shklyarov, D., \emph{Calabi-Yau structures on categories of matrix
factorizations}, J. Geom. Phys. 119 (2017), 193-207.
\item Szendroi, B., \emph{Non-commutative Donaldson-Thomas theory and the
conifold}, Geom. Topol. 12 (2008), 1171-1202.
\item Toen, B., \emph{The homotopy theory of dg-categories and derived
Morita theory}, Invent. Math. 167 (2007), 615-667.
\item Van den Bergh, M., \emph{Three-dimensional flops and
non-commutative rings}, Duke Math. J. 122 (2004), 423-455.
\end{itemize}

\end{document}
